\subsection{A nonlinear check: two dimensional Navier Stokes}\label{m:sec:expns}

The experiments above are conjugate by design. This one is not, and it is the test that matters for Sections~\ref{m:sec:loc} and~\ref{m:sec:ns}: whether the localization theory survives contact with a nonlinear forward map. We solve \eqref{m:eq:nsproj} in vorticity form on the $2\pi$ periodic torus by a pseudospectral method with two thirds dealiasing and Heun stepping, on a $24^2$ grid with $\nu=5\times10^{-3}$, $T=1$ and $80$ steps; one forward solve costs about $15$\,ms. The unknown is the initial vorticity in $D=6$ real cosine modes with wavevectors $(1,0),(0,1),(1,1),(1,-1),(2,0),(0,2)$, the prior is independent Gaussian with standard deviations $1.5(\lambda_j/\lambda_1)^{-1}$, and the truth is a prior draw. The data are $N=80$ space time points with both velocity components observed, noise $\sigma=0.45$ against a signal of standard deviation about $0.95$, and observation times uniform on $[0,T]$, so the theory applies in its $T_0=0$ form. Accuracy is the per coordinate $\Wone_1$ distance to an adaptive random walk Metropolis reference, computed as the mean absolute difference of the two quantile functions at $199$ equally spaced levels from $0.005$ to $0.995$, normalized by that coordinate's posterior standard deviation and averaged over coordinates. The reference chain runs $6000$ iterations and keeps the last $4500$, with a Gaussian proposal of the Laplace covariance whose scale is adapted every $200$ iterations toward an acceptance rate of $0.30$; its effective sample size is about $250$ per coordinate. The generators of this section use the summary \eqref{m:eq:regsummary} at the maximum a posteriori estimate, with its Jacobian by finite differences (two Gauss Newton steps in the lower panel of Table~\ref{m:tab:ns}), and resample their training residuals rather than draw the Gaussian noise of Definition~\ref{m:def:gen}. The comparator is the Laplace approximation at the maximum a posteriori estimate, which in this six mode truncation is the finite sample form of the classical parametric Bernstein von Mises limit, the finite dimensional counterpart of the limit of Konen and Nickl~\cite{konennickl2026}, and is free: one Gauss Newton solve and one Hessian.

\begin{table}[tp]\centering\footnotesize
\caption{Two dimensional Navier Stokes, $D=6$, $N=80$. Distance to the Metropolis reference in units of posterior standard deviation, per coordinate and averaged; lower is better. The upper panel varies the architecture at $M=1600$, with a one step summary and, in the localized rows, radius $c=3$; the lower panel applies the three corrections of Section~\ref{m:sec:loc} in turn, at $c=5$, $M=1100$ and $k=2$ Gauss Newton steps (Remark~\ref{m:rem:nsdiag}). Every script scores against the same reference chain and draws its own Laplace sample, so the Laplace row moves between panels; across the six runs it lies between $0.076$ and $0.083$. The last row of that panel adds to the rescaled generator a shift of the center by the cubic correction of Remark~\ref{m:rem:relin}, the third derivative term of the Laplace expansion of the posterior mean, which is not a step of Algorithm~\ref{m:alg:loc}; on the same draws that shift raises the error of the rescaled generator from $0.070$ to $0.072$, that of the centered generator from $0.091$ to $0.092$ and that of the Laplace approximation from $0.078$ to $0.081$. Over twenty further resamplings of the $8000$ draws of the rescaled generator the shift raises its error in every one of them, by $0.002$ to $0.003$. The last line is not a method: it is the Monte Carlo standard error of the reference's own posterior mean.}\label{m:tab:ns}
\begin{tabular}{@{}>{\raggedright\arraybackslash}p{3.3in}c@{}}
\toprule
& \shortstack{$\Wone_1$ /\\ posterior s.d.}\\
\midrule
\multicolumn{2}{@{}l}{\emph{Architecture, $M=1600$}}\\
\quad training from the prior, identity head & 1.077\\
\quad training from the prior, learned affine head & 0.901\\
\quad training from the prior, $K=200$ nearest summaries, affine head & 0.364\\
\quad localized training, identity head & 0.187\\
\quad localized training, learned affine head & 0.180\\
\midrule
\multicolumn{2}{@{}l}{\emph{The corrections of Section~\ref{m:sec:loc}, $c=5$, $k=2$, $M=1100$}}\\
\quad localized, uncentered, spread as measured & 0.136\\
\quad $+$ residuals centered, Theorem~\ref{m:thm:locbias}(ii) & 0.092\\
\quad $+$ spread rescaled, Theorem~\ref{m:thm:locbias}(iii) & 0.070\\
\quad $+$ center shifted by the cubic correction, Remark~\ref{m:rem:relin} & 0.072\\
\midrule
Laplace at the maximum a posteriori estimate & 0.078\\
\midrule
\emph{Monte Carlo error of the reference itself} & \emph{0.063}\\
\bottomrule
\end{tabular}
\end{table}

Table~\ref{m:tab:ns} reads as a sequence of corrections, each one predicted by Section~\ref{m:sec:loc}.

\emph{Prior training with the identity head fails, and does not improve with $M$.} The error is about one posterior standard deviation and moves from 1.131 at $M=400$ to 1.078 at $M=800$ and 1.077 at $M=1600$, so it is bias and not Monte Carlo error. This is the shrinking neighborhood caveat of Remark~\ref{m:rem:affine} made concrete: the identity head is justified on the localization set, draws from the prior lie far outside it, and there the summary \eqref{m:eq:regsummary}, whose Jacobian sits at the preliminary estimate, is not a one step efficient estimator but an arbitrary linear functional. Learning an affine head instead recovers only part of the loss, because the map is genuinely nonlinear over the prior's support.

\emph{Localization is necessary, and post hoc localization is not a substitute.} Drawing the training parameters from a proposal centered at the preliminary estimate cuts the error by a factor of 5.8 and restores the decrease in $M$. Localizing a prior trained pool after the fact by keeping the $K$ nearest summaries reaches only 0.364 at $K=200$ and degrades as $K$ grows, which is the bias variance trade off of approximate Bayesian computation and is why Algorithm~\ref{m:alg:loc} localizes the proposal rather than the pool.

Two further knobs, importance weighting the localized draws back to the prior and relinearizing the summary, are reported in Remark~\ref{m:rem:nsdiag}.

\emph{The residual mean is the one Theorem~\ref{m:thm:locbias}(ii) predicts.} The uncentered head, which resamples the training residuals, adds their sample mean to the observed summary, and so displaces its center. Per coordinate that mean is right in sign and within $0.09$ posterior standard deviations in magnitude of the linear Gaussian prediction $P\bar\theta_N$ of Theorem~\ref{m:thm:locbias}(ii), the measured mean being the smaller in five of the six coordinates (both vectors are in Section~\ref{m:sec:loc}). We attribute the difference to the nonlinearity of the map; this run does not isolate it, since the sampling error of a mean of $M=1100$ residuals is about $0.03$ posterior standard deviations and four of the six differences are no larger than that. Centering removes it, at which point the center of the GBC posterior is displaced from the reference's own posterior mean by 0.061, which is exactly the maximum a posteriori estimate's displacement. Rescaling the spread by the reciprocal of the predicted inflation $(s_j+c^2p_j^2)^{1/2}$ of Theorem~\ref{m:thm:locbias}(iii) brings the error to 0.070, which is where Algorithm~\ref{m:alg:loc} stops; shifting the center as well, by the cubic correction described in the caption of Table~\ref{m:tab:ns}, gives 0.072, so that correction does not help here. In the last coordinate at $c=5$ the theorem predicts an inflation of 1.25, and the standard deviation of the residuals is measured at 1.25 times the Laplace standard deviation, which is 1.18 times the reference's.

Two further regimes were run, each with its own reference chain of $3000$ retained draws, to look for the crossover at which non Gaussianity of the posterior would favor a simulation based method. Both used $N=25$ space time points on $100$ time steps, $M=1200$ draws from a proposal of radius $c=2$, a one step summary and the identity head. The first raises the amplitude, the constant $1.5$ in the prior standard deviations raised to $2$, which raises the proxy $\max_j|\theta_{0,j}|/\nu$ of the advection to dissipation ratio of the realized truth from $185$ at the baseline constants to $246$, still below the $347$ of the main run's truth; the second widens the posterior, with noise standard deviation $1.30$ in place of $0.45$ and that constant raised to $3$ (proxy $370$). With the residuals neither centered nor rescaled, as in the runs recorded with our paper, GBC scored 0.105 and 0.253 against the Laplace approximation's 0.074 and 0.111. With the corrections of Section~\ref{m:sec:loc} applied to the same residuals, centering and then rescaling the spread, GBC scores 0.086 and then 0.080 in the first regime and 0.172 and then 0.114 in the second, so that the corrected generator ends 0.006 and 0.003 above the Laplace approximation, while the resampling noise of a single score is 0.003 to 0.004 for each method and the Monte Carlo standard errors of the two references are 0.081 and 0.083. The crossover was not found, and neither regime made the posterior measurably less Gaussian than the main run. Their mean absolute marginal skewness, 0.10 and 0.13, varies between chain seeds about as much as between regimes (from 0.06 to 0.10 in the first, from 0.07 to 0.13 in the second and from 0.09 to 0.12 at the baseline constants), and it is of the size that sampling error alone gives a Gaussian chain of that length and autocorrelation, 0.09 on average, with 0.05 to 0.15 between the 5th and 95th percentiles; the reference chain of the main run has 0.17, against 0.07 for a Gaussian chain like it. The uncorrected generator lost more in these regimes, where the prior fractions are larger, up to 0.36 and 0.65 against 0.17 in the main run: the displacement and inflation that the two corrections remove grow with the prior fraction (Theorem~\ref{m:thm:locbias}), and once they are removed, what separates GBC from the Laplace approximation is within the noise of the comparison. Both errors are governed by the curvature of $\Gcal$ measured in posterior standard deviations, so they do not move independently and the crossover cannot be assumed.

What the table does not show is a win. Corrected, GBC reaches 0.070 against the Laplace approximation's 0.078, but over resamplings of its own draws its error averages 0.077 with standard deviation 0.005, and the reference's own Monte Carlo standard error is 0.063, so the two are indistinguishable at the resolution available: bringing that standard error down to the size of the difference in the table would need a chain about sixty times longer. In this regime the Laplace approximation is excellent and free, and recommending a simulation based method on accuracy grounds would be wrong. What the table does show is that the distance between a naive GBC implementation and a correct one is a factor of 15 (1.077 at $M=1600$ against 0.070 at $M=1100$), and that every step of that factor is prescribed by Section~\ref{m:sec:loc} rather than being a tuning choice; the cubic correction of the last row is not one of those steps. The case for GBC here is amortization, and the regimes in which it is the only option are the ones the theory excludes: $T_0>0$, model error, non Gaussian noise.

\begin{remark}[Two diagnostics from the Navier Stokes run]\label{m:rem:nsdiag}
\emph{Importance weighting back to the prior degenerates, as the algorithm says it does.} Correcting the localized proposal by weights $w\propto\pi/q$ raises the error of the learned affine head, with the one step summary and $M=1200$, at every radius at which the algorithm operates: from 0.247 to 0.279 at $c=1$, from 0.175 to 0.387 at $c=2$, from 0.172 to 0.327 at $c=3$ and from 0.269 to 0.717 at $c=5$, with effective sample sizes of 29, 6.5, 15 and 1.5 out of $M$. At $c=8$ the weighted error, 0.377 at an effective sample size of 34, is below the unweighted 0.631; the statement concerns the radii $c\le5$ at which the algorithm operates. The proposal is much narrower than the prior, so the weights grow like $e^{+\norm\theta^2}$ in its tails. The unweighted estimator is biased for the prior based posterior, because the proposal acts as a prior on the training draws; in the linear Gaussian model the bias is of first order in the prior fraction (Theorem~\ref{m:thm:locbias}) and, for the affine head, in $c^{-2}$ as well, so it is small when the data dominate the prior and the radius is not small, while the weighted estimator's variance is not.

\emph{Linearizing the summary once is what the $c^4$ term of Proposition~\ref{m:prop:radius} costs.} At $c=2$ and $M=500$ the one step and two step Gauss Newton summaries are indistinguishable for the learned affine head, 0.169 against 0.171, as they must be where the $c^4$ term is negligible. At $c=6$ the one step summary costs that head a factor of 2.3, 0.370 against 0.162, and relinearizing recovers all of it (Remark~\ref{m:rem:relin} gives the iteration and its cost).
\end{remark}

\subsection{What these experiments do not establish}\label{m:sec:expnot}

Most of them are run on a linear, one dimensional, conjugate problem. Those test the summary, the generator and the resolution accounting of Sections~\ref{m:sec:suff}, \ref{m:sec:qnn} and \ref{m:sec:resolution}, and they cannot test the nonlinear theory at all. Section~\ref{m:sec:expns} supplies the nonlinear test. It is a six dimensional Galerkin truncation, which is a small problem; the comparison is against a Metropolis reference whose own Monte Carlo error is of the same size as the differences being reported; and the truth is a prior draw, so nothing there speaks to misspecification. What it does establish is that the three corrections of Section~\ref{m:sec:loc} are the ones that matter on a genuinely nonlinear map, in the direction and roughly the magnitude the theorems predict. Whether the same holds at operational resolution, where $D$ is in the thousands and the reference is unavailable, is left open. Section~\ref{m:sec:ns} remains a specification of the pipeline at the resolution the framework is stated for, not a report of a computation at that resolution.
